No prefácio de 1859, Marx chama de base econômica o conjunto das relações de produção sobre o qual se ergue uma superestrutura jurídica e política; distingue ainda os conflitos da vida material das formas pelas quais as pessoas tomam consciência deles e os enfrentam (Marx, “Prefácio”, pp.~47--48).
Raymond Williams lê essa determinação como um jogo de pressões e limites que pesam sobre as práticas sem lhes prescrever de antemão um único desfecho (Williams, “Base e Superestrutura”, pp.~43--45).
Que o aluguel vença no fim do mês ajuda a entender por que a promessa de ganhar mais dinheiro encontra escuta; não diz, por si só, que voto sairá daquela casa.
